\section{Choose byte meanings before choosing instructions}\label{maze-p02}

The arrays have different meanings despite using the same entry width. A wall byte describes the input map. A state byte distinguishes unseen, waiting, and processed cells. A parent byte names the predecessor that first discovered a cell. A distance byte counts edges from the start. A queue byte holds a work item, and a path byte marks membership in the final route.

The sentinel 255 is possible in an unsigned byte and impossible as a cell ID in this grid. It also lies beyond every finite shortest distance: a simple route in a graph of 64 cells needs at most 63 edges. Those bounds make byte parents and distances sufficient for this assignment. A larger graph would need a fresh width and sentinel analysis.

Load IDs and sentinels with \texttt{lbu}. Loading 255 with \texttt{lb} produces minus one after sign extension, which is a different register value from unsigned 255. The direction tables are different: their negative displacements are intentional, so \texttt{lb} is the appropriate operation there.
